\documentclass{article}
% Shared layout. Keep benchmark prose and result-specific plots outside this folder.
\PassOptionsToPackage{numbers}{natbib}
\usepackage[preprint]{neurips_2026}
\usepackage[T1]{fontenc}
\usepackage[hidelinks]{hyperref}
\usepackage{booktabs,amsmath,microtype,pgfplots}
\pgfplotsset{compat=1.17}
\usetikzlibrary{patterns}
\definecolor{okblue}{HTML}{0072B2}
\definecolor{okorange}{HTML}{D55E00}
\pgfplotsset{every axis/.append style={font=\small, tick label style={font=\small}, width=\linewidth, height=5cm, grid=major}}

\newcommand{\PaperTitle}{Latency Budget: Grounded Answers Under Explicit Deadlines}
\newcommand{\PaperAbstract}{An answer can be grounded and still arrive too late to be useful. Latency Budget measures this distinction on 36 synthetic directory, calendar, and document tasks. Credit requires a single correct candidate, replayed evidence, and completion within an explicit deadline. We run 1656 local episodes across 23 repeated conditions without paid inference. The privileged oracle obtains 100.0 [90.4, 100.0] percent deadline credit (95\% Wilson interval); each blind control obtains 0.0 [0.0, 9.6] percent. An additional simulated factorial experiment varies declared thinking time and redundant tool calls while holding task deadlines fixed. At 80 ms per turn, one redundant call reduces deadline credit to 75.0 [58.9, 86.2] percent, while grounded correctness remains 100.0 [90.4, 100.0] percent. The experiment demonstrates sensitivity to wasted actions under the chosen cost model. It does not measure learned-agent performance, network latency, or human deadline preferences.}
\newcommand{\AuthorOrcid}{0009-0005-9964-6841}
\newcommand{\TaskCount}{36}
\newcommand{\ConditionCount}{23}
\newcommand{\EpisodeCount}{1656}
\newcommand{\CILevel}{95}
\newcommand{\BaseThink}{80}
\newcommand{\OracleRate}{100.0 [90.4, 100.0]\%}
\newcommand{\FloorRate}{0.0 [0.0, 9.6]\%}
\newcommand{\OracleMedian}{725.5}
\newcommand{\OracleTail}{2099}
\newcommand{\OneExtraRate}{75.0 [58.9, 86.2]\%}
\newcommand{\TwoExtraRate}{30.6 [18.0, 46.9]\%}
\newcommand{\FourExtraRate}{0.0 [0.0, 9.6]\%}
\newcommand{\SlowRate}{27.8 [15.8, 44.0]\%}
\newcommand{\AllCorrectRate}{100.0 [90.4, 100.0]\%}
\newcommand{\OneExtraLate}{25.0 [13.8, 41.1]\%}
\newcommand{\SlowThink}{320}
\newcommand{\DoubleThink}{160}
\newcommand{\ZeroThink}{0}
\newcommand{\ExtraCount}{4}
\newcommand{\RepeatCount}{2}
\newcommand{\ControlCount}{7}
\newcommand{\GridCount}{16}
\newcommand{\BootDraws}{2000}
\newcommand{\SweepSeed}{20261029}
\newcommand{\SourceCommit}{c82bd8290465}
\newcommand{\NullDelta}{0.0 [0.0, 0.0]}
\newcommand{\OneDelta}{25.0 [16.7, 33.3]}

\title{\PaperTitle}
\author{Pooria Arab\\Independent researcher\\ORCID \texttt{\AuthorOrcid}}
\begin{document}
\maketitle
\begin{abstract}
\PaperAbstract
\end{abstract}
\section{Why this matters}
An assistant can find the right information and still miss the moment when a user needs it. A directory lookup during conversation differs from a background research job. Both require correct evidence, but their time limits differ. An evaluation that records only answer accuracy treats these situations as equivalent. An evaluation that records only speed rewards immediate but unsupported replies. Neither measure alone captures useful completion under a deadline.\par
Latency Budget defines a narrower target: produce a grounded answer before an explicit time limit. Grounded means that a replayed successful tool result contains the expected value. Timely means that model time and scheduled tool costs remain within the task budget. These requirements are conjunctive. Correctness cannot cancel lateness, and spare time cannot cancel an incorrect answer. The benchmark also reports each component so a failed headline score remains interpretable.\par
This paper studies that measurement rule through controlled execution. The main contribution is a benchmark design with an explicit time model, replayed evidence, and separate late-correct reporting. A further contribution is a simulated sensitivity experiment. It changes wasted tool actions and thinking time while preserving the answer and task deadline. This isolates time accounting from answer generation. The study does not establish the relative capability of language models.\par
The distinction matters for interpretation. A benchmark may separate deliberately different policies without representing the full range of real agents. Our oracle has access to the reference path and expected answer. It verifies attainability, not reasoning ability. Our blind controls check whether easy shortcuts receive credit. The simulated sweep checks whether the measurement responds to its intended inputs. Each experiment supports a different part of the validity argument.\par
\section{Related work}
HELM treats efficiency as one evaluation dimension alongside other model properties~\cite{liang2022helm}. That broad approach motivates reporting more than answer correctness. Latency Budget instead selects a small interactive setting where time changes whether an answer earns credit. It does not replace broad language evaluation. It makes the deadline rule explicit enough to inspect through task traces.\par
Kapoor et al. argue that agent evaluation should account for resources alongside accuracy~\cite{kapoor2024agents}. Our design follows that principle for elapsed budgeted time. Time and financial cost remain distinct: a free script can waste a deadline, while an expensive model can finish promptly. The present study incurs no inference charges, but this says nothing about the deployment cost of a future learned agent.\par
The Agentic Benchmark Checklist emphasizes task setup, outcome validity, and reporting~\cite{zhu2025abc}. We use privileged and trivial controls to inspect the outcome rule. A full score for the privileged control establishes a reachable positive case. Failed blind controls establish limited negative cases. They do not constitute a general proof against reward exploitation. New answer styles or different tool sequences may expose failures outside these probes.\par
Miller discusses uncertainty and paired comparisons in language-model evaluation~\cite{miller2024errors}. We retain the task as the analysis unit and preserve task pairing in controller comparisons. We also distinguish empirical uncertainty from actual sampling. The public tasks are a designed suite. Their intervals describe a hypothetical sampling interpretation, rather than variation from a measured user population.\par
The family audit describes Latency Budget as a synthetic deadline benchmark with declared or measured billing~\cite{arab2026family}. We retain that scope and study the declared mode in depth. PayBench provides a useful construct-first presentation pattern~\cite{plunkett2026paybench}. Its payment decisions, preference labels, and matched safe/unsafe pairs do not transfer to this task. Here the natural comparison is the same grounded answer under different time costs.\par
\section{Experiment design}
\subsection{Task, observation, and outcome}
Each episode starts with a fictional user request, a stated deadline, and a small tool set. Tasks ask for a directory value, a calendar room, or information connected through documents and people. The agent chooses calls and submits a final answer. The benchmark observes the tool sequence and final output. The outcome combines candidate selection, evidence recovery, and time accounting.\par
The task supplies a finite list of possible answer values to the evaluator. The final answer must name exactly one candidate, and that candidate must equal the expected value. Matching ignores letter case and permits surrounding words, subject to word boundaries. This is a lexical rule. It does not establish that the surrounding prose affirms the candidate or expresses a valid explanation. The intended interface therefore requests the value alone.\par
Grounding uses replay rather than trusting submitted tool results. The evaluator reconstructs tool outcomes against the task fixtures. A successful result must contain the expected value. A call whose own arguments already contain that value cannot establish evidence. This prevents a simple echo from appearing to verify an answer. Tool descriptions expose cost information but do not count as answer evidence.\par
The environment includes directory lookup, calendar retrieval, document search, and cost discovery. Directory lookup returns a person's room and extension. Calendar retrieval returns events for a date. Document search returns matching records. These services are deterministic for the selected fixture. The study therefore isolates evaluation behavior without failures from a live external service.\par
\subsection{Reference paths and budget classes}
The public suite contains \TaskCount{} tasks. The classes are named voice, UI, and background. Table~\ref{tab:coverage} gives their counts, reference path lengths, and deadline ranges. These names describe intended budget regimes. They do not certify suitability for a spoken interface, screen interaction, or background production system. No user study sets these thresholds.\par
\begin{table}[htbp]
\centering\small
\caption{Task coverage. Calls refer to the supplied reference path. Deadline values use virtual milliseconds.}\label{tab:coverage}
\begin{tabular}{lrrrr}\toprule
Class & Tasks & Calls & Minimum deadline & Maximum deadline \\\midrule
voice & 12 & 1 & 447 & 641 \\
ui & 12 & 2 & 875 & 1494 \\
background & 12 & 3 & 2025 & 3714 \\\bottomrule
\end{tabular}
\end{table}

A reference path lists the calls used to reach the expected value. Voice tasks use direct retrieval. UI tasks link a calendar or document record to a directory entry. Background tasks follow longer links among those sources. The budget derives from the supplied reference path under the baseline declared thinking time. This gives the positive control enough room to finish while making additional work potentially costly.\par
Reference-path length and budget class change together. A class comparison therefore cannot isolate time pressure from reasoning depth, tool composition, or fixture structure. A lower voice score would not alone establish that an agent is worse at conversational tasks. A stronger causal design would hold task content fixed while assigning multiple deadline regimes. The simulated sweep makes a narrower within-task comparison.\par
\begin{figure}[htbp]\centering
\begin{tikzpicture}\begin{axis}[ylabel={Virtual milliseconds},xlabel={Task index: voice, UI, then background},legend pos=north west]
\addplot+[mark=*,mark size=1pt] table[x=x,y=time,col sep=comma]{figures/coverage.generated.csv};\addlegendentry{Oracle time}
\addplot+[mark=*,mark size=1pt] table[x=x,y=deadline,col sep=comma]{figures/coverage.generated.csv};\addlegendentry{Deadline}
\end{axis}\end{tikzpicture}
\caption{Reference paths fit their assigned budgets. The same ordering shows longer paths and looser deadlines together; class effects are confounded with path length.}\label{fig:coverage}\end{figure}

\subsection{Clock model and experimental boundaries}
All reported experiments use declared time. For each recorded call, the evaluator charges its scheduled tool latency and a declared thinking cost. It also charges thinking time for the final answer. There is no sleep that reproduces these virtual costs. Actual script execution time does not determine the result. The model describes sequential work under a fixed schedule, rather than network behavior.\par
The implementation also has a measured mode. It records wall-clock gaps around model and environment activity while retaining scheduled tool costs. That mode answers a different question and would require hardware, transport, and load controls. We do not combine its measurements with declared-mode results. A future comparison must report which part of its time comes from observation and which part comes from a schedule.\par
The declared mode bills a flat tool log as separate turns. It cannot infer concurrency from an unordered list of calls. As a result, batching does not earn the latency reduction of real parallel execution. This conservative rule is transparent but incomplete. Agents that depend on parallel services need a task model with explicit dependencies and overlapping intervals.\par
\section{Coverage and failure taxonomy}
Coverage centers on retrieval with short dependencies. The fixtures contain fictional people, room names, events, and document records. Candidate values include short identifiers and document titles. This range exercises different lexical forms while keeping grading deterministic. It excludes open-ended synthesis, partial answers, speech production, uncertain evidence, and sustained conversation.\par
The public generator and answers make this a transparent instrument, not a secure evaluation set. A policy can memorize public values. Replaying evidence makes a bare memorized answer insufficient, but it does not prevent a memorized tool plan. The privileged oracle intentionally has this access. Its result must never be presented as evidence of generalization.\par
The outcome taxonomy separates four broad states: incorrect or missing answer, correct text without accepted evidence, grounded but late answer, and grounded answer within budget. The last state earns credit. The taxonomy is operational rather than psychological. It records what the evaluator can establish, not why the agent chose an action.\par
A missing answer includes an empty output or a stopped episode with no final value. An incorrect answer includes a wrong candidate, no recognized candidate, or several named candidates. These cases need different model repairs, but all fail the correctness condition. A refusal can be fast and still fail an answerable retrieval task. This makes fast abstention an unsuitable way to maximize the headline score.\par
An ungrounded answer is a distinct failure even when its text matches the key. It could reflect guessing, memory, or an inadequate trace. Replay establishes whether accepted evidence exists under the environment rules. It does not establish that the policy used that evidence in its reasoning. The evaluator observes available support, not an internal causal explanation.\par
A late correct answer preserves both text correctness and accepted evidence, but exceeds the time limit. This state is the main diagnostic for this study. Without a separate count, it would be indistinguishable from an incorrect answer in the headline mean. A deadline curve shows whether additional time would recover credit for the same trace.\par
A further process taxonomy can identify redundant calls, costly discovery, repeated failures, or excessive model time. Our simulated study manipulates only redundant successful calls and declared thinking time. It does not test error recovery or discovery strategy. Those mechanisms can share a final outcome while requiring different interventions.\par
\section{Metrics and scoring}
\subsection{Conjunctive deadline credit}
Let $A_i$ indicate that task $i$ has exactly one correct named candidate with accepted replayed evidence. Let $D_i$ be its deadline. A trace has $m_i$ calls with scheduled costs $\ell_{ij}$. Under declared thinking time $h$, its budgeted completion time is\par
\begin{equation} T_i(h)=h(m_i+1)+\sum_{j=1}^{m_i}\ell_{ij}. \end{equation}
The extra thinking charge represents the final answer. Deadline credit is $Q_i=A_i\,\mathbf{1}\{T_i\leq D_i\}$. The suite score is the arithmetic mean of $Q_i$. A late correct indicator is $L_i=A_i\,\mathbf{1}\{T_i>D_i\}$. Thus $A_i=Q_i+L_i$ for every complete timed trace. We use this identity as an analysis check.\par
This formulation exposes a practical tradeoff. Additional evidence collection may improve correctness but consume time. The simulated experiment holds correctness fixed, so it measures only the time side. It cannot estimate how a learned agent's extra reasoning changes answer quality. Removing all extra computation may damage a real agent even when it helps the privileged scripted policy.\par
\subsection{Time summaries and deadline curves}
We report the median and upper-tail percentile of virtual time over all submitted traces in a condition. These summaries include incorrect answers. Reporting only successful completion times would hide slow failures and favor policies that give up. The percentile uses interpolation between ordered observations. It is descriptive; this paper does not attach population precision to its tail estimate.\par
For a multiplier $f$, the deadline curve rescales the original deadline and scores the same answer and calls. Each point is the mean of $A_i\mathbf{1}\{T_i\leq\operatorname{round}(fD_i)\}$. A curve cannot fall as the multiplier increases. Its maximum cannot exceed grounded correctness. These properties give useful consistency checks without assuming that more time changes the agent's behavior.\par
A fixed-trace curve is not an adaptive policy evaluation. The agent does not see the changed budget and choose a new path. The curve asks whether the existing trace would fit a different acceptance threshold. This distinction prevents a common interpretation error: a higher score under relaxed grading does not prove that an agent uses extra time productively.\par
\subsection{Uncertainty and comparisons}
Every reported binary rate includes a \CILevel{}\% Wilson interval. The denominator is the number of tasks in that condition or class. Repeated deterministic runs confirm reproducibility and do not create additional independent tasks. The same rule applies when the simulated grid contains many conditions. We never pool repeated cells into a larger apparent sample of agent capabilities.\par
Wilson intervals avoid a zero-width normal interval at a perfect or empty score. They still require a hypothetical sampling model. These tasks share generation templates and fixed class proportions. Correlation and limited template diversity can make interval coverage differ from that assumption. The finite-suite counts remain exact, regardless of whether the interval interpretation fits a deployment population.\par
Controller comparisons pair each task with itself. We subtract the score with redundant calls from the score without those calls at the same thinking time. Bootstrap resampling preserves the class counts and task pairing. We use \BootDraws{} draws with seed \SweepSeed{}. The interval measures variation across tasks under this resampling rule. It is not uncertainty about simulator correctness.\par
When every paired difference is identical, the bootstrap interval collapses. That is a boundary property of the empirical distribution, not proof of universal effectiveness. We show component Wilson intervals beside these comparisons. We make no significance claims, no multiple-model ranking, and no estimate of human benefit.\par
\section{Results}
\subsection{Scripted positive and negative controls}
The preregistered study completes \EpisodeCount{} episodes across \ConditionCount{} conditions. Every condition runs \RepeatCount{} times on the full task set. The baseline controls use declared thinking time \BaseThink{} ms. There are \ControlCount{} baseline policies and \GridCount{} simulated grid cells. All inference is local scripted execution.\par
The packaged positive policy follows the supplied reference path and submits the expected answer. It earns \OracleRate{} deadline credit. Its median virtual time is \OracleMedian{} ms, and its p95 is \OracleTail{} ms. These results show that the fixtures and chosen budgets admit the intended reference behavior. They do not test whether an unassisted model can discover that behavior.\par
Each negative control earns \FloorRate{} deadline credit. The packaged bad policy says yes. The constant policy says OK. Empty, echo, and refusal policies submit their designated text. The first-tool policy calls directory lookup with empty arguments and then submits an empty answer. This limited first-tool definition is fixed before collection; it does not search for valid arguments or act as a retrieval heuristic.\par
\begin{table}[htbp]
\centering\small
\caption{Scripted controls. Rates are percentages with task-level Wilson intervals. Speed alone earns no credit.}\label{tab:controls}
\begin{tabular}{lrrrr}\toprule
Policy & Deadline credit & Correct & Late correct & p50 ms \\\midrule
packaged-good & 100.0 [90.4, 100.0] & 100.0 [90.4, 100.0] & 0.0 [0.0, 9.6] & 725.5 \\
packaged-bad & 0.0 [0.0, 9.6] & 0.0 [0.0, 9.6] & 0.0 [0.0, 9.6] & 80 \\
empty & 0.0 [0.0, 9.6] & 0.0 [0.0, 9.6] & 0.0 [0.0, 9.6] & 80 \\
constant & 0.0 [0.0, 9.6] & 0.0 [0.0, 9.6] & 0.0 [0.0, 9.6] & 80 \\
echo & 0.0 [0.0, 9.6] & 0.0 [0.0, 9.6] & 0.0 [0.0, 9.6] & 80 \\
refuse & 0.0 [0.0, 9.6] & 0.0 [0.0, 9.6] & 0.0 [0.0, 9.6] & 80 \\
first-tool & 0.0 [0.0, 9.6] & 0.0 [0.0, 9.6] & 0.0 [0.0, 9.6] & 401 \\\bottomrule
\end{tabular}
\end{table}

The fast controls demonstrate why speed cannot stand alone. Their short traces consume little time, but none supply a grounded correct answer. Conversely, the oracle establishes a positive case under the same scoring rule. Together these checks exclude a constant scorer on the tested inputs. They do not exclude more targeted strategies that exploit lexical or fixture weaknesses.\par
\begin{figure}[htbp]\centering
\begin{tikzpicture}\begin{axis}[ymin=0,ymax=105,ylabel={Deadline credit (percent)},xtick={1,2,3,4,5,6,7},xticklabels={Oracle,Bad,Empty,Constant,Echo,Refuse,First tool}]
\addplot+[okblue,only marks,mark=square*,error bars/.cd,y dir=both,y explicit] table[x=x,y=y,y error minus=lo,y error plus=hi,col sep=comma]{figures/controls.generated.csv};
\end{axis}\end{tikzpicture}
\caption{Only the privileged oracle earns deadline credit. Whiskers show Wilson intervals across tasks. The controls use declared virtual time.}\label{fig:controls}\end{figure}

\subsection{Simulated factorial experiment}
The additional simulated experiment changes thinking time and redundant call count. Each simulated policy first completes the privileged reference path. It then repeats the final successful call before giving the expected answer. Each repeat returns the same evidence. The task deadline and tool schedule stay fixed. Only declared thinking time and the number of calls change.\par
All simulated cells retain grounded correctness of \AllCorrectRate{}. At \BaseThink{} ms per turn, a simulated extra call reduces deadline credit to \OneExtraRate{}. The corresponding late-correct rate is \OneExtraLate{}. With a simulated pair of extra calls, credit is \TwoExtraRate{}. With \ExtraCount{} simulated extra calls, it is \FourExtraRate{}. The decomposition attributes these score changes to lateness under the specified model.\par
The effect of simulated thinking time is also conditional on slack. At \DoubleThink{} ms and no redundant calls, the simulated policy still reaches the baseline ceiling shown in Table~\ref{tab:grid}. At \SlowThink{} ms without redundant calls, its credit falls to \SlowRate{}. A larger thinking charge can therefore exhaust the budget even when the tool plan stays minimal. These are simulated thresholds, not observed model speeds.\par
\begin{table}[htbp]
\centering\small
\caption{Simulated factorial experiment. All cells retain the same grounded correctness shown in the text. Rate intervals use the original tasks, not repeated episodes.}\label{tab:grid}
\begin{tabular}{rrrr}\toprule
Think ms & Extra calls & Deadline credit & Late correct \\\midrule
0 & 0 & 100.0 [90.4, 100.0] & 0.0 [0.0, 9.6] \\
0 & 1 & 100.0 [90.4, 100.0] & 0.0 [0.0, 9.6] \\
0 & 2 & 86.1 [71.3, 93.9] & 13.9 [6.1, 28.7] \\
0 & 4 & 13.9 [6.1, 28.7] & 86.1 [71.3, 93.9] \\
80 & 0 & 100.0 [90.4, 100.0] & 0.0 [0.0, 9.6] \\
80 & 1 & 75.0 [58.9, 86.2] & 25.0 [13.8, 41.1] \\
80 & 2 & 30.6 [18.0, 46.9] & 69.4 [53.1, 82.0] \\
80 & 4 & 0.0 [0.0, 9.6] & 100.0 [90.4, 100.0] \\
160 & 0 & 100.0 [90.4, 100.0] & 0.0 [0.0, 9.6] \\
160 & 1 & 38.9 [24.8, 55.1] & 61.1 [44.9, 75.2] \\
160 & 2 & 2.8 [0.5, 14.2] & 97.2 [85.8, 99.5] \\
160 & 4 & 0.0 [0.0, 9.6] & 100.0 [90.4, 100.0] \\
320 & 0 & 27.8 [15.8, 44.0] & 72.2 [56.0, 84.2] \\
320 & 1 & 0.0 [0.0, 9.6] & 100.0 [90.4, 100.0] \\
320 & 2 & 0.0 [0.0, 9.6] & 100.0 [90.4, 100.0] \\
320 & 4 & 0.0 [0.0, 9.6] & 100.0 [90.4, 100.0] \\\bottomrule
\end{tabular}
\end{table}

\begin{figure}[htbp]\centering
\begin{tikzpicture}\begin{axis}[ymin=0,ymax=105,ylabel={Deadline credit (percent)},xlabel={Simulated thinking time per turn (ms)},xtick={0,80,160,320},legend style={at={(0.5,-0.3)},anchor=north,legend columns=2}]
\addplot+[okblue,mark=*,error bars/.cd,y dir=both,y explicit] table[x=x,y=y,y error minus=lo,y error plus=hi,col sep=comma]{figures/sweep0.generated.csv};\addlegendentry{0 extra calls}
\addplot+[okorange,mark=square*,error bars/.cd,y dir=both,y explicit] table[x=x,y=y,y error minus=lo,y error plus=hi,col sep=comma]{figures/sweep1.generated.csv};\addlegendentry{1 extra calls}
\addplot+[black,mark=triangle*,error bars/.cd,y dir=both,y explicit] table[x=x,y=y,y error minus=lo,y error plus=hi,col sep=comma]{figures/sweep2.generated.csv};\addlegendentry{2 extra calls}
\addplot+[gray,mark=diamond*,error bars/.cd,y dir=both,y explicit] table[x=x,y=y,y error minus=lo,y error plus=hi,col sep=comma]{figures/sweep4.generated.csv};\addlegendentry{4 extra calls}
\end{axis}\end{tikzpicture}
\caption{Simulated thinking time and redundant calls jointly consume the fixed budget. Whiskers show task-level Wilson intervals; connected points are not a fitted model.}\label{fig:sweep}\end{figure}

The simulated response surface follows a simple mechanism. Adding a call contributes its scheduled latency and another thinking charge. Increasing thinking time affects every call and the final answer. The effects therefore interact through the number of recorded actions. This interaction is built into the cost model. The experiment checks that executed traces and scores exhibit it; it is not a newly discovered law of agent behavior.\par
\subsection{Simulated duplicate-removal control}
We treat the simulated zero-extra-call condition as an ideal duplicate-removal controller. It keeps the reference path and removes only the deliberately inserted repeats. It knows which calls are redundant by construction. This is an ablation of avoidable work. It does not evaluate a detector that must distinguish useful retries from duplicates in an unfamiliar task.\par
At the baseline thinking time, removing a simulated extra call improves deadline credit by \OneDelta{} percentage points. At zero simulated thinking time, removing that call changes credit by \NullDelta{} percentage points. The latter is a null result. A time saving need not change a binary score when both paths remain inside the budget.\par
\begin{table}[htbp]
\centering\small
\caption{Simulated duplicate-removal control. Paired, class-stratified bootstrap intervals accompany each improvement. Boundary intervals can collapse.}\label{tab:contrasts}
\begin{tabular}{rrr}\toprule
Think ms & Removed calls & Improvement, percentage points \\\midrule
0 & 1 & 0.0 [0.0, 0.0] \\
0 & 2 & 13.9 [5.6, 25.0] \\
0 & 4 & 86.1 [75.0, 97.2] \\
80 & 1 & 25.0 [16.7, 33.3] \\
80 & 2 & 69.4 [55.6, 83.3] \\
80 & 4 & 100.0 [100.0, 100.0] \\
160 & 1 & 61.1 [50.0, 69.4] \\
160 & 2 & 97.2 [91.7, 100.0] \\
160 & 4 & 100.0 [100.0, 100.0] \\
320 & 1 & 27.8 [19.4, 33.3] \\
320 & 2 & 27.8 [19.4, 33.3] \\
320 & 4 & 27.8 [19.4, 33.3] \\\bottomrule
\end{tabular}
\end{table}

At the slowest simulated thinking setting, duplicate removal cannot restore full deadline credit. The minimal trace itself exceeds some budgets. This bounds the intervention: eliminating wasted actions helps only the part of latency that those actions cause. It cannot remove the time charged for required work. A practical controller must distinguish avoidable delay from a fundamentally mismatched deadline.\par
\subsection{Simulated deadline relaxation and repeatability}
Figure~\ref{fig:curve} rescales deadlines for simulated traces at the baseline thinking setting. Each point uses the same evidence and final answer. Relaxation recovers credit that the original budget denies. The curves therefore reveal a gap between grounded correctness and timely usefulness without changing the task solution.\par
\begin{figure}[htbp]\centering
\begin{tikzpicture}\begin{axis}[ymin=0,ymax=105,ylabel={Deadline credit (percent)},xlabel={Deadline multiplier},xtick={0.5,1,2,4},legend style={at={(0.5,-0.3)},anchor=north,legend columns=2}]
\addplot+[okblue,mark=*,error bars/.cd,y dir=both,y explicit] table[x=x,y=y,y error minus=lo,y error plus=hi,col sep=comma]{figures/curve0.generated.csv};\addlegendentry{0 extra calls}
\addplot+[okorange,mark=square*,error bars/.cd,y dir=both,y explicit] table[x=x,y=y,y error minus=lo,y error plus=hi,col sep=comma]{figures/curve1.generated.csv};\addlegendentry{1 extra calls}
\addplot+[black,mark=triangle*,error bars/.cd,y dir=both,y explicit] table[x=x,y=y,y error minus=lo,y error plus=hi,col sep=comma]{figures/curve2.generated.csv};\addlegendentry{2 extra calls}
\addplot+[gray,mark=diamond*,error bars/.cd,y dir=both,y explicit] table[x=x,y=y,y error minus=lo,y error plus=hi,col sep=comma]{figures/curve4.generated.csv};\addlegendentry{4 extra calls}
\end{axis}\end{tikzpicture}
\caption{Simulated deadline relaxation restores credit to late correct traces. Thinking time stays at the baseline setting. Curves rescore fixed behavior; they do not measure adaptation. Wilson intervals show task uncertainty.}\label{fig:curve}\end{figure}

Every repeated condition produces the same semantic trace hash. The comparison retains task identifiers, scores, reasons, answers, tool logs, transcripts, and clock mode. It excludes incidental wall-clock stamps because declared time does not use them. The raw repeats remain available for inspection. This is reproducibility evidence for the local scripted protocol, not a claim that future hosted models will behave deterministically.\par
The headline results support the registered directional checks. The oracle reaches its intended ceiling, tested blind policies fail, and simulated extra work cannot increase credit. The simulated duplicate-removal comparison also includes a documented null. No task is dropped based on its score, and no learned-model result is inferred from these controls.\par
\section{Roadmap}
The next study should replace privileged scripts with agents that receive only the ordinary prompt and tool observations. It should freeze model identifiers, decoding settings, tool interfaces, and call limits before collection. A free-only access guard can enforce the inference budget. Provider availability and publication terms need verification at that time. No hosted-model availability probe forms part of the present study.\par
A second priority is to separate deadline class from task depth. The same information need should appear under several independently assigned budgets. This would support a causal estimate of deadline pressure on policy behavior. It would also distinguish failure to solve a task from failure to finish it quickly. Class labels alone cannot provide that estimate in the current suite.\par
A richer time model should represent service variance, queueing, cancellation, and overlapping tool calls. These changes require explicit dependency graphs and a stated clock trust boundary. A fixed schedule remains useful for controlled comparisons, but it should not impersonate a production trace. A future measurement study should calibrate virtual costs against observed services and publish the calibration errors.\par
Finally, deadline selection needs validation against actual use. Users may accept a partial answer early and a full answer later. Some tasks tolerate delay but penalize interruptions. Others require an acknowledgement before the final result. A single terminal deadline intentionally omits these utility functions. Human evaluation should precede claims that the chosen thresholds optimize user experience.\par
\section{Caveats and validity limits}
\paragraph{Synthetic task distribution.} Fictional retrieval tasks reduce ambiguity and simplify replay. They also omit tool failures, data changes, access control, and ambiguous user goals. The suite is small and public. Shared templates constrain diversity even when names or values change. A new seed supplies different instances, not an independent task domain.\par
\paragraph{Privileged policies.} The oracle and simulated controls read the supplied path and expected answer. This is deliberate for instrument testing. It leaks the solution by design and cannot support a model leaderboard. The blind controls avoid that privileged information, but they are weak policies rather than plausible task solvers. No learned agent or human participant is evaluated.\par
\paragraph{Lexical grading.} Candidate extraction checks text, not sentence meaning. A negated sentence can still name a candidate. Evidence checks whether a replayed result contains the expected value, not whether the full chain logically establishes it. These limits matter when moving beyond value-only replies. This paper does not claim a complete adversarial audit of the grader.\par
\paragraph{Path optimality.} The supplied reference path is a construction target. It is not a proof that no shorter grounded trace exists. Some downstream arguments are available to a privileged policy before it follows the intended chain. A future audit should search for alternative traces and distinguish necessary dependency from narrative intent. Deadline calibration should then state whether it uses an authored path or a verified optimum.\par
\paragraph{Clock trust.} Declared time is an assumption, not an observed service measurement. Measured records need a trusted environment; a client can otherwise claim unrealistically short intervals. Local and remote scoring interfaces may treat clock records differently. All results here avoid that comparison by fixing declared mode. No production latency guarantee follows from this study.\par
\paragraph{Interval scope.} Wilson intervals describe hypothetical task sampling. The stratified bootstrap describes empirical paired variation. Neither captures uncertainty about whether the benchmark represents useful work. Repeating deterministic traces confirms implementation stability but does not narrow those validity gaps. A large number of episodes should not obscure the small set of distinct tasks.\par
\paragraph{Unrun work.} No free hosted model is selected or contacted. Availability elsewhere is unknown. There is no human deadline study, independent grader audit, cross-domain replication, or external preregistration deposit. The plan is committed before data collection within the study history. That establishes ordering in the artifact, not an independent registration guarantee.\par
\paragraph{Responsible use.} The fixtures contain fictional records and do not request private user data. The benchmark should support diagnosis of deadline behavior, not certification of safety-critical response time. Optimizing a hard deadline alone may encourage shallow answers outside the evaluated tasks. Any deployment decision also needs evidence about correctness, failure handling, and the costs of delay in its own setting.\par
\section{Citation and availability}
The paper and generated figures use CC-BY-4.0. Benchmark code and task data retain their MIT license. The reproducible study packet, raw traces, source hashes, and citation metadata are available at \url{https://github.com/pooriaarab/benchmark-latency-budget}. The frozen benchmark source is \texttt{\SourceCommit}. The packet supplies the exact repeat command and distinguishes generated data from authored analysis. No venue publication is implied.\par
\section{Authors}
Pooria Arab is an independent researcher. The author ORCID appears above. AI assistance supports study implementation, analysis checks, and drafting. There is no claim of independent review or human participant validation. Responsibility for the study's claims and release remains with the author.\par
\bibliographystyle{plainnat}
\bibliography{references}
\end{document}
